\documentclass[11pt,twoside]{article}

\usepackage[margin=1in]{geometry}
\usepackage{amsmath,amssymb,amsthm,mathtools}
\usepackage{array}
\usepackage[inline,shortlabels]{enumitem}
\usepackage{xcolor}
\usepackage{microtype}
\usepackage{../exercises/course-exercises}

% Revision date: maintained by scripts/build_lecture_notes.py. Keep this command.
\newcommand{\lecturerevised}{2026-10-25 21:25 UTC-06:00}
\usepackage{eso-pic}
\AddToShipoutPictureBG{%
  \AtPageLowerLeft{%
    \put(\LenToUnit{0.5\paperwidth},\LenToUnit{0.5in}){%
      \makebox(0,0){\normalfont\footnotesize\color{black!60}Last revised: \lecturerevised}%
    }%
  }%
}

\usepackage[square,authoryear]{natbib}
\usepackage{xr-hyper}
\usepackage[colorlinks=true,linkcolor=blue!60!black,
  urlcolor=blue!60!black,citecolor=blue!60!black]{hyperref}
\externaldocument[lecture2-]{lecture02}
\hypersetup{
  pdftitle={Lecture 13: Representation learning and transfer},
  pdfauthor={Csaba Szepesvari}
}
\usepackage[nameinlink,capitalize]{cleveref}
\crefname{exercise}{Exercise}{Exercises}
\Crefname{exercise}{Exercise}{Exercises}

\usepackage{enotez}
\DeclareInstance{enotez-list}{lecture}{list}{format=\normalsize,number={#1.}}
\setenotez{list-name=Endnotes,list-style=lecture,backref=true}
\crefname{endnote}{Endnote}{Endnotes}

\definecolor{courseblue}{HTML}{17365D}
\newcommand{\E}{\mathbb E}
\newcommand{\I}{\mathbb I}
\renewcommand{\P}{\mathbb P}
\newcommand{\R}{\mathbb R}
\newcommand{\KL}{D}
\newcommand{\TV}{\mathrm{TV}}
\makeatletter
\@for\@letter:=A,B,C,D,E,F,G,H,I,J,K,L,M,N,O,P,Q,R,S,T,U,V,W,X,Y,Z\do{%
  \expandafter\edef\csname c\@letter\endcsname{\noexpand\mathcal{\@letter}}}
\makeatother
\newcommand{\lecturenumber}{13}

\newtheorem{theorem}{Theorem}[section]
\newtheorem{proposition}[theorem]{Proposition}
\newtheorem{corollary}[theorem]{Corollary}
\newtheorem{lemma}[theorem]{Lemma}
\theoremstyle{definition}
\newtheorem{definition}[theorem]{Definition}
\newtheorem{example}[theorem]{Example}
\theoremstyle{remark}
\newtheorem{remark}[theorem]{Remark}

\renewcommand{\thepage}{13--\arabic{page}}
\renewcommand{\thesection}{13.\arabic{section}}
\renewcommand{\theequation}{13.\arabic{equation}}
\pagestyle{myheadings}
\markboth{Lecture 13: Representation learning and transfer}
{Lecture 13: Representation learning and transfer}

\begin{document}

\begin{center}
\fcolorbox{courseblue}{white}{%
\begin{minipage}{0.92\textwidth}
\vspace{1mm}
\textbf{CMPUT 654: Mathematical Foundations of Modern AI Systems}
\hfill Fall 2026

\vspace{3mm}
\begin{center}
{\Large\bfseries Lecture 13: Representation learning and transfer}

\vspace{1mm}
October 15, 2026
\end{center}

\vspace{2mm}
\textit{Lecturer: Csaba Szepesv\'ari}
\hfill
\textit{Note prepared by: Csaba Szepesv\'ari}
\vspace{1mm}
\end{minipage}}
\end{center}

\noindent\textbf{Reading guide.}
Neural networks, including language models, are often said to be powerful
because they learn representations. This is usually contrasted with
``linear'' methods where a feature map needs to be hand-picked
for the method to work well.

We first define representation-learning algorithms in the context of learning
a map and connect the definition to neural networks. Representation learning
has a clear formulation in the multitask setting: learn one feature map from
several tasks, with a separate predictor for each task. We ask whether sharing
reduces the samples needed to predict on those tasks. Transfer learning asks a
further question: do the learned features help a new task? We give a
finite-class analysis of each problem, then an efficient transfer algorithm
for a linear model. The statistical and computational questions are separate:
how much data suffices, and how can the learner compute useful features?

The main results are the joint-ERM bounds in
\cref{thm:consistent-multitask,thm:agnostic-multitask}, the separation from
independent learning in \cref{prop:boxes-sharing}, and the transfer guarantees
in \cref{thm:calibrated-transfer,cor:linear-transfer}. The proofs use the
finite-class argument from Lecture 2 and concentration from Lecture 6.
The examples also connect to parameter sharing in Lecture 4 and coverage in
Lecture 7.

\section{Representation learning}
\label{sec:representation-learning}

Let $\cX$ be an input space, $\cY$ a label space, and $\cZ$ a feature space.
We consider learning a predictor $\cX\to\cY$ from $n$ labeled examples.

\begin{definition}[Representation-learning algorithm]
\label{def:representation-learning}
Let $\cH$ be a class of maps $\cX\to\cZ$ and $\cG$ a class of maps
$\cZ\to\cY$. A \emph{representation-learning algorithm}
for sample size $n$ with representation
class $\cH$ and head class $\cG$ is a map
\[
 \cA:(\cX\times\cY)^n\longrightarrow\cH\times\cG.
\]
\end{definition}


Writing $\cA(S)=(h_S,g_S)$, the learner returns a \emph{representation} (map)
$h_S$ and a \emph{head} $g_S$. The resulting predictor is
$\widehat f_S=g_S\circ h_S$. For example, $h_S$ can compute a neural
network's hidden features and $g_S$ can classify using those features.
As can be seen from the definition, the main distinctive feature of a representation
learner is the signature of the underlying map. Namely,
that it returns a map that is explicitly decomposed into two maps. This definition
is meant to capture the essence of these methods and as such is a bit oversimplified.
Every supervised learning algorithm is a trivial representation learner:
Just choose $\cZ = \cX$, $\cH$ consist of the $\cX \to \cX$ identity map.
Of course,
we are more interested in nontrivial representation learners, for which for at least
some sample $S\in (\cX\times \cY)^n$, $h_S$ is not the identity map (if $\cZ = \cX$).

Our primary interest in representation learning is to gain insight
into why learning with neural networks is special as compared to simpler alternatives.
To explain this connection recall that a feedforward neural network is
given by a directed (computation) graph, with nodes performing simple computation
over the information sent to them along the edges ending at them,
where learning changes the actual computation performed.
Now, take any minimum-size
cut that separates the inputs (sources of the graph) from the outputs (sinks of the graph):
This is a set of nodes such that removing this set, no input is connected to no output,
while no subset of this set of nodes has the same property.
If the nodes form a layered graph (with respect to distance to inputs), as
it happens for example, in multi-layered networks, any node set in a layer
will be such cut. The map determined by the part of the computation graph up to the cut
defines $h:x\mapsto z$; the remaining computation defines
$g:z\mapsto\widehat y$, while the whole network computes $g\circ h$.

It follows that neural networks are in a way ``super representation learners'':
They are considering many compositional maps simultaneously (corresponding to all possible
cuts) and I suspect this is one of their strengths. However, for the sake
of this lecture we stick to a simpler view that ignores this multi-representation
learning aspect, focusing on just a one-way decomposition of the target
function.

% TODO: This is at a bad spot; at one point we want to write about
% how within a single task we often have multiple ``implicit'' tasks.
% For example, in a language model, the input sequences may come from
% different languages, natural or artificial. Here, learning each
% language can be considered its own task. But what can go even deeper
% and argue that in language modeling you need to learn grammar,
% and semantics, which are different tasks. It is is not entirely
% clear though how this maps to our simple definition of representation
% learning.
%
% Shared features also occur within a single prediction problem. In a
% multiclass classifier, class scores $s_c=g_c(h(x))$ use a common
% representation. Softmax cross-entropy couples the scores through the loss
% \[
%  \ell(s,y)=-s_y+\log\sum_c e^{s_c}.
% \]
% The class-specific score functions are trained through this joint loss.
% In autoregressive language modeling, the same network computes features
% across contexts and sequence positions, as in Lectures 1 and 4.

\section{Multitask learning}
\label{sec:multitask-learning}
The usefulness of representation learning is easiest to argue for
when the learner needs to solve multiple prediction problems,
that share a common representation. In this section
we study this problem
in \emph{multitask learning},
where a learner needs to find predictors for several tasks and
is evaluated on those tasks.
We extend the learner's output to one
representation and $T$ heads. The shared-map, separate-head structure makes
it easy to articulate  the potential benefits of learning a shared representation.

\subsection{The learning problem}
\label{sec:setup}

The learner receives labeled data from each task of the $t\in [T]$ tasks
and returns a predictor for
each task. For the first statistical comparison, take $\cH$ and $\cG$ to be
finite and nonempty. The finite-class model isolates the accounting created
by a shared representation: the representation complexity is paid once
across tasks, while the head complexity is paid once per task. More general
representation-class bounds replace logarithmic cardinalities by other
complexity measures and retain this separation.
Such class-level bounds characterize the statistical behavior of the class.
Explaining which representation neural-network training selects from a large
architecture class additionally requires properties of the training procedure
and the learned representation. The same selection problem is central to
neural-network generalization theory beyond multitask and representation
learning.
In the Gaussian linear model studied later,
we analyze an explicit representation-learning strategy
and identify what an algorithm-level guarantee requires.

Let  $P_1,\ldots,P_T$ on $\cX\times\cY$ be the data-generating laws of the
individual $T$ tasks.
For each task $t\in [T]$, the sample is
\[
 S_t=((X_{t1},Y_{t1}),\ldots,(X_{tn},Y_{tn}))\sim P_t^{\otimes n}.
\]
All $nT$ pairs are independent.
The task identity of each example is given.\endnote{
We choose the multi-task setting of \cref{sec:setup}  where the task
identities are known and different heads needs to be learnt for each task
for pedagogical reasons.
An alternative, and in many ways more common scenario, is
when the data is not grouped based on tasks (thus the
information of which task a given example belongs to is hidden)
and/or the same
head is to be used across all the tasks.
Examples abound; one obvious case is language modeling
with multiple languages (e.g., English, German, Hungarian, Chinese, Python, Java, $\dots$).
In these domain-agnostic (or task-agnostic) setting,
the representation learnt could separate the language agnostic
content from the language identity,
thus making the learning of the head across many different
languages more efficient.
}
The laws may differ, including their input marginals.
No statistical model over the tasks is assumed.

Let $\ell:\cY\times\cY\to[0,1]$ be a loss. Write
\begin{equation}
 R_t(f)=\int\ell(f(x),y)\,P_t(dx,dy),
 \qquad
 R_{S_t}(f)=\frac1n\sum_{j=1}^n\ell(f(X_{tj}),Y_{tj}).
 \label{eq:task-risks}
\end{equation}
A sample in the subscript marks an empirical risk: $R_t$ is computed from the
law $P_t$, while $R_{S_t}$ is computed from the sample $S_t$. Write
$S=(S_1,\ldots,S_T)$ for the entire source sample.
For an output $(h,\boldsymbol g)\in\cH\times\cG^T$, with
$\boldsymbol g=(g_1,\ldots,g_T)$, the task-$t$ predictor is $g_t\circ h$.
We evaluate the output by its average population risk and use the
corresponding empirical risk for training:
\begin{equation}
 \overline R(h,\boldsymbol g)=\frac1T\sum_{t=1}^T R_t(g_t\circ h),
 \qquad
 \overline R_S(h,\boldsymbol g)
 =\frac1T\sum_{t=1}^T R_{S_t}(g_t\circ h).
 \label{eq:average-risks}
\end{equation}
A bar marks an average over the $T$ source tasks, and a hat marks a quantity
returned by a learner, such as the pair $(\widehat h,\widehat{\boldsymbol g})$
defined next.
Joint empirical risk minimization chooses one representation and all the heads:
\begin{equation}
 (\widehat h,\widehat{\boldsymbol g})
 \in\mathop{\rm argmin}_{h\in\cH,\ \boldsymbol g\in\cG^T}
       \overline R_S(h,\boldsymbol g).
 \label{eq:joint-erm}
\end{equation}
Independent learning instead chooses a separate $(h_t,g_t)$ for every task.

Whether joint learning offers a statistical advantage over independent
learning depends on the head class \(\mathcal G\).
An unrestricted head could undo
the statistical restriction imposed by $h$. The intended regime has a rich
representation class and a relatively simple head class:
$\ln|\cH|$ is large compared with $\ln|\cG|$.

\subsection{The statistical benefit of sharing}
\label{sec:joint-statistics}

We first compare joint and separate learning when one representation permits
perfect prediction on all tasks. We then allow an approximation cost for
requiring a common representation.

\subsubsection{Realizable bounded losses}

Suppose one representation $h^\star\in\cH$ and heads
$g_1^\star,\ldots,g_T^\star\in\cG$ achieve zero risk on every source task.
Because the loss is nonnegative, every source loss of this predictor tuple
equals zero almost surely. Joint ERM therefore has zero empirical loss almost
surely. The consistency argument from
\cref{lecture2-thm:finite-class-consistency} extends as follows to the bounded
loss from \cref{sec:setup} and to independent observations whose laws may
differ.

Let $\cF$ be a finite nonempty class of predictors from $\cU$ to $\cY$.
Let $(U_1,V_1),\ldots,(U_m,V_m)$ be independent, not necessarily identically
distributed, and let $\ell:\cY\times\cY\to[0,1]$. For $f\in\cF$, define
\[
 R_i(f)=\E[\ell(f(U_i),V_i)],
 \qquad
 \overline R_m(f)=\frac1m\sum_{i=1}^m R_i(f).
\]
Suppose some $f^\star\in\cF$ satisfies $R_i(f^\star)=0$ for every $i\in[m]$.
On the same experiment, let $\widehat f$ be an $\cF$-valued random predictor
with
\[
 \ell(\widehat f(U_i),V_i)=0
 \qquad\text{for every }i\in[m]
\]
almost surely. Such a predictor exists because $R_i(f^\star)=0$ and the loss
is nonnegative.

\begin{proposition}[Finite-class zero-loss bound]
\label{prop:independent-finite-class-consistency}
For every $\delta\in(0,1)$,
\[
 \P\left\{
 \overline R_m(\widehat f)
 \le \frac{\ln|\cF|+\ln(1/\delta)}{m}
 \right\}\ge1-\delta.
\]
Moreover,
\[
 \E[\overline R_m(\widehat f)]
 \le \frac{\ln|\cF|+1}{m}.
\]
\end{proposition}

\Cref{ex:l13-bounded-loss-consistency} asks you to prove the proposition. The
key inequality is
\[
 \P\bigl(\ell(f(U_i),V_i)=0\bigr)\le 1-R_i(f),
\]
which follows from $0\le\ell\le1$. For $0$--$1$ loss, the inequality is an
equality and the proposition reduces to the usual consistency bound.

To apply the proposition, augment each input with its task identity and
consider the class
\[
 \cF_{\rm joint}
 =\{(t,x)\mapsto g_t(h(x)):
       h\in\cH,\ (g_1,\ldots,g_T)\in\cG^T\}.
\]
The class has cardinality at most $|\cH||\cG|^T$. For the $nT$ augmented
observations $((t,X_{tj}),Y_{tj})$, the quantity $\overline R_{nT}$ in
\cref{prop:independent-finite-class-consistency} equals
$\overline R(h,\boldsymbol g)$.
The subscripts $\mathrm{src}$ and $\mathrm{trg}$ used below abbreviate source
and target.
For $\delta\in(0,1)$, put
\[
 a_{\rm src}(\delta)
 =\frac{\ln|\cH|+T\ln|\cG|+\ln(1/\delta)}{nT}.
\]

\begin{theorem}[Joint zero-loss bound]
\label{thm:consistent-multitask}
Under the zero-risk source assumption above, let
$(\widehat h,\widehat{\boldsymbol g})$ be any predictor tuple with zero loss
on every source example. For $0<\delta<1$,
\begin{equation}
 \P\left\{
 \overline R(\widehat h,\widehat{\boldsymbol g})
 \le a_{\rm src}(\delta)
 \right\}\ge1-\delta.
 \label{eq:consistent-multitask}
\end{equation}
The expected average risk also satisfies
\begin{equation}
 \E[\overline R(\widehat h,\widehat{\boldsymbol g})]
 \le \frac{\ln|\cH|+T\ln|\cG|+1}{nT}.
 \label{eq:expected-multitask}
\end{equation}
\end{theorem}

\begin{proof}
Apply \cref{prop:independent-finite-class-consistency} with
$m=nT$ and $|\cF_{\rm joint}|\le|\cH||\cG|^T$.
\end{proof}

The bound separates the two statistical costs:
\[
 a_{\rm src}(\delta)
 =\underbrace{\frac{\ln|\cH|}{nT}}_{\text{one shared representation}}
  +\underbrace{\frac{\ln|\cG|}{n}}_{\text{one head per task}}
  +\frac{\ln(1/\delta)}{nT}.
\]
For the comparison below, write
\[
 \overline R_{\rm shared}
 :=\overline R(\widehat h,\widehat{\boldsymbol g})
\]
for the risk of the shared learner's output. Independent learning returns a
separate pair $(\widehat h_t^{\rm ind},\widehat g_t^{\rm ind})\in
\cH\times\cG$ for each task $t$; define its average risk by
\[
 \overline R_{\rm independent}
 :=\frac1T\sum_{t=1}^T
 R_t(\widehat g_t^{\rm ind}\circ\widehat h_t^{\rm ind}).
\]
The class of independent predictor tuples has size at most
$(|\cH||\cG|)^T$. The same finite-class result gives the sufficient bound
\begin{equation}
 \P\left\{
 \overline R_{\rm independent}
 \le \frac{\ln|\cH|+\ln|\cG|}{n}
       +\frac{\ln(1/\delta)}{nT}
 \right\}\ge1-\delta.
 \label{eq:independent-upper}
\end{equation}
Thus the representation term in the shared upper bound is smaller by a
factor of $T$. Comparing upper bounds is a useful first comparison. The next
example also gives a lower bound for independent learning.

Lecture 4 analyzed pooling across positions when parameters are shared.
Here the pooled object is a feature map rather than a single probability.
In both cases the statistical gain depends on a restriction that is correct
for the tasks being pooled.

\subsubsection{An example where sharing provably helps}
\label{sec:boxes-sharing}

Return to Lecture 2's boxes. Let $\cX=[d]$ with $d\ge2$. An unknown map
$h^\star:[d]\to\{0,1\}$ assigns a bit to every box.
Task $t$ has an unknown bit $b_t$ and label
\begin{equation}
 Y=h^\star(X)\mathbin{\oplus} b_t,
 \qquad X\text{ uniform on }[d],
 \label{eq:label-reversal-model}
\end{equation}
where $\oplus$ denotes addition modulo two. All training inputs are
independent. Choose all $2^d$ Boolean maps as $\cH$ and the two heads
$g_b(z)=z\oplus b$ as $\cG$.

\begin{proposition}[Shared boxes versus independent learning]
\label{prop:boxes-sharing}
Joint ERM in this model has
\begin{equation}
 \E[\overline R_{\rm shared}]
 \le \frac{(d+T)\ln2+1}{nT}.
 \label{eq:boxes-shared-upper}
\end{equation}
For any collection of learners whose task-$t$ predictor uses only $S_t$
and independent randomization, there is a fixed choice of
$(h^\star,b_1,\ldots,b_T)$ for which
\begin{equation}
 \E[\overline R_{\rm independent}]
 \ge \frac12\left(1-\frac1d\right)^n.
 \label{eq:boxes-independent-lower}
\end{equation}
\end{proposition}

\begin{proof}
The upper bound is \cref{eq:expected-multitask}.
For the lower bound, set every $b_t=0$ and temporarily draw the $d$ values of
$h^\star$ independently as fair bits, independently of all inputs and learner
randomization. A fresh uniform test box is absent from $S_t$ with
probability $(1-1/d)^n$. Conditional on that event and the task-$t$ data,
its bit remains fair. Every task-$t$ prediction therefore has conditional
error $1/2$ on the event. Average over tasks and then over the random map.
At least one fixed map has average expected error at least the displayed
lower bound.
\end{proof}

For a fixed error tolerance $\varepsilon>0$, the shared upper bound shows
that
\[
 n\ge
 \left\lceil\frac{(d+T)\ln2+1}{\varepsilon T}\right\rceil
\]
examples per task suffice to make
$\E[\overline R_{\rm shared}]\le\varepsilon$. Thus
$O(1+d/T)$ examples per task suffice for any fixed tolerance. In contrast,
independent learning needs $n=\Omega(d)$ to guarantee a fixed small tolerance
uniformly over this family.
For example, Bernoulli's inequality gives an independent-learning error of
at least $1/4$ whenever $n\le d/2$. When $T$ is of order $d$, the shared
upper bound can be made smaller than $1/4$ with a constant $n$.
These are statistical statements; enumerating all representations need not
be computationally efficient.

\subsubsection{Misspecification and joint ERM}

We now allow arbitrary task laws and any loss in $[0,1]$. Define
\begin{equation}
 R_t^\star=\min_{h\in\cH,\ g\in\cG}R_t(g\circ h),
 \qquad
 A_t(h)=\min_{g\in\cG}R_t(g\circ h)-R_t^\star.
 \label{eq:representation-error}
\end{equation}
The quantity $A_t(h)\ge0$ is the extra risk imposed by using representation
$h$, even with its best head. The benchmark $R_t^\star$ allows task $t$ its
own representation. It need not equal Bayes risk: the composite class itself
may be misspecified.

Write
\begin{equation}
 \overline A(h)=\frac1T\sum_{t=1}^T A_t(h),
 \qquad
 \Delta_{\rm share}=\min_{h\in\cH}\overline A(h).
 \label{eq:sharing-cost}
\end{equation}
The number $\Delta_{\rm share}$ is the approximation cost of requiring one
representation. Because the heads can be chosen separately,
\begin{equation}
 \min_{h,\boldsymbol g}\overline R(h,\boldsymbol g)
 =\frac1T\sum_{t=1}^T R_t^\star+\Delta_{\rm share}.
 \label{eq:sharing-benchmark}
\end{equation}
Exact sharing means $\Delta_{\rm share}=0$, not necessarily zero label noise.

Allow the source learner an optimization error $\eta_{\rm src}\ge0$:
\begin{equation}
 \overline R_S(\widehat h,\widehat{\boldsymbol g})
 \le\min_{h,\boldsymbol g}\overline R_S(h,\boldsymbol g)
       +\eta_{\rm src}.
 \label{eq:approximate-joint-erm}
\end{equation}

\begin{theorem}[Agnostic joint ERM]
\label{thm:agnostic-multitask}
Under \cref{sec:setup} and \cref{eq:approximate-joint-erm}, define
\[
 b_{\rm src}(\delta)=
 \sqrt{\frac{2\{\ln|\cH|+T\ln|\cG|+\ln(2/\delta)\}}{nT}}.
\]
With probability at least $1-\delta$,
\begin{align}
 \overline R(\widehat h,\widehat{\boldsymbol g})
   -\frac1T\sum_{t=1}^T R_t^\star
 &\le \Delta_{\rm share}+\eta_{\rm src}+b_{\rm src}(\delta),
 \label{eq:agnostic-joint}\\
 \overline A(\widehat h)
 &\le \Delta_{\rm share}+\eta_{\rm src}+b_{\rm src}(\delta).
 \label{eq:source-representation-bound}
\end{align}
\end{theorem}

\begin{proof}
For a fixed tuple, Hoeffding's inequality applies to the $nT$ independent,
bounded losses, although their laws differ across tasks. A union bound
over at most $K=|\cH||\cG|^T$ tuples gives
\[
 \sup_{h,\boldsymbol g}
 |\overline R_S(h,\boldsymbol g)-\overline R(h,\boldsymbol g)|
 \le e_{\rm src}:=\sqrt{\frac{\ln(2K/\delta)}{2nT}}
\]
with probability at least $1-\delta$. On this event, for a population
minimizer $(h^\circ,\boldsymbol g^\circ)$,
\begin{align*}
 \overline R(\widehat h,\widehat{\boldsymbol g})
 &\le\overline R_S(\widehat h,\widehat{\boldsymbol g})+e_{\rm src}\\
 &\le\overline R_S(h^\circ,\boldsymbol g^\circ)
       +\eta_{\rm src}+e_{\rm src}\\
 &\le\overline R(h^\circ,\boldsymbol g^\circ)
       +\eta_{\rm src}+2e_{\rm src}.
\end{align*}
Use \cref{eq:sharing-benchmark} and
$2e_{\rm src}=b_{\rm src}(\delta)$ for the first
claim. For the second, minimize each population head risk at $\widehat h$:
$\overline A(\widehat h)\le
\overline R(\widehat h,\widehat{\boldsymbol g})-T^{-1}\sum_tR_t^\star$.
\end{proof}

The estimation term has order
$\sqrt{\ln|\cH|/(nT)+\ln|\cG|/n}$.
Independent ERM has an analogous bound with order
$\sqrt{(\ln|\cH|+\ln|\cG|)/n}$ and no $\Delta_{\rm share}$.
Sharing reduces the estimation cost but can increase approximation error.
When $\Delta_{\rm share}>0$, even the best shared predictor has greater
average risk than the best separate predictors. Whether sharing hurts at a
finite sample size depends on the estimation errors as well.
The added approximation error is one source of \emph{negative transfer}:
sharing can increase
risk relative to learning separately.

\section{Transfer learning}
\label{sec:transfer}

In \emph{transfer learning}, information learned from source tasks is used to
learn a predictor for a target task. We study a two-stage procedure: estimate
$\widehat h$ from the source samples $S_1,\ldots,S_T$, then freeze
$\widehat h$ and fit a head using examples from a new target task. Learning
how to fit new tasks from experience on source tasks is also called
\emph{learning-to-learn} or \emph{meta-learning}.

Fix a target law $P_0$ and define $R_0$, $R_0^\star$, and $A_0(h)$ as in
\cref{eq:task-risks,eq:representation-error}. After estimating $\widehat h$,
the learner receives
$S_0=((X_{01},Y_{01}),\ldots,(X_{0m},Y_{0m}))\sim P_0^{\otimes m}$,
independently of all source data and source-learner randomization.
The output $(\widehat h,\widehat g_0)$ predicts by
$\widehat f_0=\widehat g_0\circ\widehat h$. The objective is small target
excess risk $R_0(\widehat f_0)-R_0^\star$.

\subsection{Finite-class analysis}

We retain the finite classes and bounded loss from \cref{sec:setup}.
\Cref{thm:agnostic-multitask} bounds $\overline A(\widehat h)$. A guarantee for a
particular task requires the source tasks to constrain the features needed
by that task. The existence of one representation that works for all tasks
leaves the behavior of the learned representation undetermined.

The source tasks constrain a candidate representation through its average
source representation error. Define the \emph{transfer envelope}
\begin{equation}
 \Psi(u)
 =
 \max\bigl(\{A_0(h):h\in\cH,\ \overline A(h)\le u\}\cup\{0\}\bigr),
 \qquad u\ge0.
 \label{eq:transfer-envelope}
\end{equation}
The value $\Psi(u)$ is the largest target representation error among the
candidate representations whose average source representation error is at
most $u$. Thus $\Psi$ states exactly what source performance permits on the
target. The envelope is nondecreasing and, by construction,
\begin{equation}
 A_0(h)\le\Psi(\overline A(h))
 \qquad\text{for every }h\in\cH.
 \label{eq:transfer-calibration}
\end{equation}

\begin{proposition}[Sharpness of the transfer envelope]
\label{prop:transfer-envelope-sharpness}
Let $\phi:[0,\infty)\to[0,\infty)$ be nondecreasing and suppose
\[
 A_0(h)\le\phi(\overline A(h))
 \qquad\text{for every }h\in\cH.
\]
Then $\Psi(u)\le\phi(u)$ for every $u\ge0$. Moreover, whenever
$\{h\in\cH:\overline A(h)\le u\}$ is nonempty, some $h_u\in\cH$ satisfies
\[
 \overline A(h_u)\le u,
 \qquad
 A_0(h_u)=\Psi(u).
\]
\end{proposition}

\Cref{ex:l13-transfer-envelope} asks you to prove the proposition and apply
it to the dropped-feature example below. Thus $\Psi(u)$ is the exact worst
case over candidate representations compatible with the source-error
constraint $\overline A(h)\le u$. A maximizing representation witnesses that
the source-error constraint permits target representation error $\Psi(u)$.
An algorithm-specific guarantee can be smaller by using information beyond
the value of $\overline A(\widehat h)$.

A useful transfer guarantee proves that $\Psi(u)$ is small in the relevant
range. A common bound is
\begin{equation}
 \Psi(u)\le\kappa u+\beta,
 \qquad \kappa,\beta\ge0.
 \label{eq:linear-calibration}
\end{equation}
The multiplier measures how strongly source accuracy constrains target
accuracy. The additive term allows imperfect transfer.
The examples in \cref{sec:calibration-examples} bound the envelope by
comparing the errors induced by every candidate representation.

\subsubsection{A source-to-target theorem}

With $\widehat h$ frozen, target training searches only over the head class.
Choose a head with
\begin{equation}
 R_{S_0}(\widehat g_0\circ\widehat h)
 \le\min_{g\in\cG}R_{S_0}(g\circ\widehat h)+\eta_{\rm trg},
 \label{eq:target-erm}
\end{equation}
where $\eta_{\rm trg}\ge0$ and $R_{S_0}$ uses the $m$ target examples.

\begin{theorem}[Transfer from source representation error]
\label{thm:calibrated-transfer}
Assume \cref{eq:target-erm}. Suppose that, for some
$\varepsilon_{\rm src}\ge0$ and $0<\delta_{\rm src},\delta_{\rm trg}<1$,
the source procedure satisfies
\[
 \P\{\overline A(\widehat h)\le\varepsilon_{\rm src}\}
 \ge1-\delta_{\rm src}.
\]
Put
\[
 b_{\rm trg}(\delta_{\rm trg})
 =\sqrt{\frac{2\ln(2|\cG|/\delta_{\rm trg})}m}.
\]
With probability at least $1-\delta_{\rm src}-\delta_{\rm trg}$,
\begin{equation}
 R_0(\widehat g_0\circ\widehat h)-R_0^\star
 \le\Psi(\varepsilon_{\rm src})
     +\eta_{\rm trg}+b_{\rm trg}(\delta_{\rm trg}).
 \label{eq:general-transfer}
\end{equation}
In particular, agnostic joint ERM gives
\begin{equation}
 R_0(\widehat g_0\circ\widehat h)-R_0^\star
 \le
 \Psi\!\left(
   \Delta_{\rm share}+\eta_{\rm src}
   +b_{\rm src}(\delta_{\rm src})\right)
 +\eta_{\rm trg}+b_{\rm trg}(\delta_{\rm trg}).
 \label{eq:erm-transfer}
\end{equation}
For an exactly optimized, jointly realizable source problem, one may
instead take
$\varepsilon_{\rm src}=a_{\rm src}(\delta_{\rm src})$.
\end{theorem}

\begin{proof}
Condition on the source data and source-learner randomization, so that
$\widehat h$ is fixed and the target
sample still has law $P_0^{\otimes m}$. Apply the Hoeffding--union-bound
argument to the $|\cG|$ predictors $g\circ\widehat h$. With conditional
probability at least $1-\delta_{\rm trg}$,
\[
 R_0(\widehat g_0\circ\widehat h)
 -\min_{g\in\cG}R_0(g\circ\widehat h)
 \le\eta_{\rm trg}+b_{\rm trg}(\delta_{\rm trg}).
\]
The conditional failure probability is at most $\delta_{\rm trg}$ for every source
realization, so the unconditional failure probability is also at most
$\delta_{\rm trg}$. On the intersection with the source event,
\[
 \min_gR_0(g\circ\widehat h)-R_0^\star
 =A_0(\widehat h)
 \le\Psi(\overline A(\widehat h))
 \le\Psi(\varepsilon_{\rm src}).
\]
Adding the inequalities proves \cref{eq:general-transfer}.
Use \cref{thm:agnostic-multitask} or
\cref{thm:consistent-multitask} for the respective choices of
$\varepsilon_{\rm src}$. In the realizable case,
$\overline A(\widehat h)\le
\overline R(\widehat h,\widehat{\boldsymbol g})$.
\end{proof}

The target fitting cost contains $\ln|\cG|$, not $\ln|\cH|$.
Learning the target from scratch has the finite-class ERM bound
\[
 R_0(\widehat f_{\rm scratch})-R_0^\star
 \le \sqrt{
 \frac{2\{\ln|\cH|+\ln|\cG|+\ln(2/\delta_{\rm trg})\}}m}
\]
for an exact ERM. Transfer can improve the guarantee when
$\Psi(\varepsilon_{\rm src})$ is smaller than the estimation cost saved on the
target. The target need not be sampled from a distribution over tasks;
\cref{eq:transfer-calibration} is the explicit connection to the fixed
source tasks.

The theorem also answers the individual-source-task question. Fix
$t\in[T]$ and define
\[
 \Psi_t(u)
 =\max\bigl(\{A_t(h):h\in\cH,\ \overline A(h)\le u\}\cup\{0\}\bigr).
\]
After freezing $\widehat h$ and refitting task $t$'s head on fresh examples,
\cref{eq:general-transfer} holds for that task with $\Psi_t$.
For $q$ designated tasks, the same source event and head failure
probabilities $\delta_j$ give simultaneous guarantees with failure
probability at most $\delta_{\rm src}+\sum_{j=1}^q\delta_j$.
The guarantee uses refitting samples that are independent of the source
examples used to learn $\widehat h$. It concerns the newly fitted heads;
guarantees for the original fitted heads require corresponding
individual-task source bounds.

\subsubsection{Bounding the transfer envelope}
\label{sec:calibration-examples}

The next two examples identify source tasks that constrain target-relevant
features and source tasks that leave a necessary feature unconstrained.

\paragraph{Label-reversed tasks.}
Let's return to the example of \cref{sec:boxes-sharing}.
In \cref{eq:label-reversal-model}, let the target also have a uniform input
with the same law $\mu$ as the source tasks, and label $h^\star(X)\oplus b_0$.
For a candidate $h$, put
\[
 q(h)=\mu\bigl(\{x:h(x)\ne h^\star(x)\}\bigr).
\]
Recall that $\cG=\{g_0,g_1\}$ with $g_b(z)=z\oplus b$ for $b\in\{0,1\}$.
For each task $t\in[T]$, the target predictor is $g_{b_t}\circ h^\star$.
Since each $g_b$ is a bijection on $\{0,1\}$,
\[
 R_t(g_{b_t}\circ h)
 =\mu\bigl(\{x:g_{b_t}(h(x))\ne g_{b_t}(h^\star(x))\}\bigr)
 =q(h),
 \qquad
 R_t(g_{1-b_t}\circ h)=1-q(h).
\]
Since $R_t^* = 0$,
\begin{equation}
 A_t(h)=\min_{g\in \cG} R_t(g \circ h) =
 \min\{q(h),1-q(h)\}=A_0(h)=\overline A(h)\,.
 \label{eq:boxes-calibration}
\end{equation}
Thus $\Psi(u)\le u$, independently of $T$. The transfer theorem (\cref{thm:calibrated-transfer})) gives, for
exact source and target ERM,
\begin{equation}
 R_0(\widehat g_0\circ\widehat h)
 \le
 \frac{(d+T)\ln2+\ln(1/\delta_{\rm src})}{nT}
 +\sqrt{\frac{2\ln(4/\delta_{\rm trg})}m}.
 \label{eq:boxes-transfer}
\end{equation}
For fixed accuracy and confidence, $T$ of order $d$ permits constant $n$
and constant $m$. A learner with only the $m$ target observations still has
the expected-error lower bound $\frac12(1-1/d)^m$ for some target map.
For comparison in expectation, a risk bound of $\varepsilon$ with failure
probability $\delta$ implies expected risk at most $\varepsilon+\delta$.
Choosing both constants small gives a genuine individual-task advantage.

\paragraph{A target feature unconstrained by the source tasks.}
Let $\cX=\cZ=\{0,1\}^2$, let $\cY=\{0,1\}$, and use $0$--$1$ loss.
Let $U,V$ be independent fair bits and set $X=(U,V)$. The representation
and head classes are
\[
 \cH=\{h_{\rm full},h_{\rm drop}\},
 \qquad
 h_{\rm full}(u,v)=(u,v),
 \qquad
 h_{\rm drop}(u,v)=(u,0),
\]
and
\[
 \cG=\{\pi_1,\pi_2\},
 \qquad
 \pi_1(z_1,z_2)=z_1,
 \qquad
 \pi_2(z_1,z_2)=z_2.
\]
An observation from each source task is
$(X,Y)=((U,V),U)$, whereas a target observation is
$(X,Y)=((U,V),V)$. Thus the source and target input laws are the same and
uniform on $\{0,1\}^2$, while the label functions differ. The full
representation has zero risk on every source task using $\pi_1$ and zero
target risk using $\pi_2$, so $R_t^\star=R_0^\star=0$. For the dropped
representation,
\[
 R_t(\pi_1\circ h_{\rm drop})=0
 \quad(t\in[T]),
 \qquad
 R_0(\pi_1\circ h_{\rm drop})
 =R_0(\pi_2\circ h_{\rm drop})=\frac12.
\]
Consequently,
\[
 \overline A(h_{\rm drop})=0,
 \qquad
 A_0(h_{\rm drop})=\frac12.
\]
No source loss depends on whether a representation retains $V$. The transfer
envelope therefore satisfies $\Psi(0)=1/2$, and repeating the source tasks
leaves the target feature $V$ unconstrained.

The last two examples concern whether the source task objectives constrain the
features needed by the target. In the label-reversal example, every source
task supplies such a constraint. In the present example, no source task
constrains retention of $V$.

Lecture 7's coverage condition addresses a separate issue: comparison of
errors under different input laws. The two forms of coverage can be imposed
together. For an intermediate representation error $A_{\rm mid}(h)$, suppose
an input-law comparison gives
$A_0(h)\le\kappa_1 A_{\rm mid}(h)+\beta_1$ and a source-task comparison gives
$A_{\rm mid}(h)\le\kappa_2\overline A(h)+\beta_2$. Substitution gives
\[
 A_0(h)
 \le\kappa_1\kappa_2\overline A(h)
   +\kappa_1\beta_2+\beta_1.
\]
Thus distributional coverage and source-task coverage address different
obstacles and combine directly.

\subsection{Transfer with linear representations}
\label{sec:linear-computation}

We now give an efficient algorithm for the transfer problem. The learner
estimates a linear representation from source tasks by an eigendecomposition,
then fits a linear head on target data. The model lets us verify the
source-to-target relation and bound the risk of the resulting target
predictor.

This is the same representation-learning template as in the earlier finite-class
analysis, specialized to linear maps and squared loss. Concretely,
$\mathcal X=\mathbb R^d$, $\mathcal Z=\mathbb R^k$, $\mathcal Y=\mathbb R$, and
we use the squared loss $\ell(y,y')=(y-y')^2$. The representation and head
classes are
\begin{align}
 \cH
 &=\{h_C:\R^d\to\R^k:
       h_C(x)=C^\top x,\ C\in\R^{d\times k},\ C^\top C=I_k\},
 \label{eq:linear-representation-class}\\
 \cG
 &=\{g_b:\R^k\to\R:
       g_b(z)=b^\top z,\ b\in\R^k\}.
 \label{eq:linear-head-class}
\end{align}
The composite predictor is
$f_{C,b}(x)=g_b(h_C(x))=b^\top C^\top x=x^\top Cb$.
Thus the linear model is the previous class-based setup with explicit,
infinite classes $\cH$ and $\cG$.

\subsubsection{The learning problem}

There are $T$ source tasks and one target task, indexed by $0$.
Let $1\le k<d$. The unknown shared object is a $k$-dimensional subspace of
$\R^d$. Choose a matrix $B\in\R^{d\times k}$ whose orthonormal columns form
a basis for this subspace, and let $a_0,a_1,\ldots,a_T\in\R^k$ be the task
coordinates. The orthonormal convention preserves every full-column-rank
linear model: if a feature matrix factors as $M=BR$, replacing a task
coefficient $c_t$ by $a_t=Rc_t$ gives $Mc_t=Ba_t$. An observation from
task $t$ satisfies
\begin{equation}
 Y=X^\top Ba_t+\varepsilon,
 \qquad X\sim\cN(0,I_d),
 \label{eq:source-model}
\end{equation}
where the noise has mean zero and is independent of $X$, with finite
variance $\sigma^2$ on each source task and $\sigma_0^2$ on the target task.
The learner is given $k$, samples $S_1,\ldots,S_T$ of $n$ examples each,
and a target sample $S_0$ of $m$ examples. All covariates and noises in these
samples are mutually independent. Write $N=nT$ for the source sample size.

Thus each $C$ selects a candidate $k$-dimensional feature subspace and an
orthonormal basis for it. Because $\cG$ contains every linear head, a change
of basis changes the head coordinates but not the available composite
predictors. The normalization also makes $CC^\top$ the orthogonal projector
onto the candidate subspace, as used below.
Use $R_t$, $R_t^\star$, and $A_t$ for these classes and squared loss.
The learner estimates $\widehat B$ from the source samples alone and fits
$\widehat a_0$ on $S_0$, returning $(h_{\widehat B},g_{\widehat a_0})$.
The predictor is $\widehat f_0(x)=x^\top\widehat\theta$, where
$\widehat\theta=\widehat B\widehat a_0$. The objective is small expected
target excess risk, which in this model is
\begin{equation}
 \E[R_0(\widehat f_0)-R_0^\star]
 =\E[\|\widehat\theta-Ba_0\|_2^2],
 \qquad R_0^\star=\sigma_0^2.
 \label{eq:linear-target-objective}
\end{equation}

The statistical point is the same as in the finite-class transfer analysis: source
training constrains the representation only insofar as it rules out bad shared
subspaces. In this linear model the constraint is explicit. For a candidate
representation $C$, the squared norm of the component of $Ba_0$ orthogonal to
$\operatorname{span}(C)$ is the target representation error $A_0(C)$. The comparison is
therefore between the target representation error $A_0(C)$ and the average source
representation error $\overline A(C)$, and the source-task matrix
$D=\frac1T\sum_{t=1}^Ta_ta_t^\top$ measures how strongly the source tasks cover the
relevant directions. For well-spread sources, $D\succ0$ and a target direction with
small source representation error inherits a comparable target representation
error. This is the linear analogue of the earlier calibration bound, with the
quantitative factor $\kappa_0=a_0^\top D^{-1}a_0$.

The model gives two distinct comparisons with the earlier theory. The
statistical comparison concerns the number of target examples needed after the
source tasks have identified a useful representation. The corollary below
bounds the expected target excess risk by a source term of order
$\|a_0\|_2^2V/(nT\lambda^2)$ and a target head-fitting term
$\sigma_0^2 k/(m-k-1)$. Target-only least squares instead has target term
$\sigma_0^2 d/(m-d-1)$. Thus the ideal sharing benefit is a reduction from
ambient dimension $d$ to representation dimension $k$ in the target-label cost;
the source cost is amortized across the $T$ source tasks. The displayed bound
also quantifies the price of imperfect subspace recovery. This conclusion is a
property of the stated Gaussian linear model.

The algorithmic comparison asks how to obtain the shared representation. Joint
source ERM is a natural benchmark:
\begin{equation*}
 \min_{C^\top C=I_k,\,b_1,\ldots,b_T}
 \frac1N\sum_{t=1}^T\sum_{j=1}^n
 \bigl(Y_{tj}-b_t^\top C^\top X_{tj}\bigr)^2.
\end{equation*}
This objective is precisely the source-training objective for a two-layer
linear network with a shared first layer $C^\top$ and task-specific output
weights $b_t$. It is a finite-dimensional rank-constrained least-squares
problem. At the population level, minimizing over the heads makes its excess
objective
\[
 \operatorname{tr}\!\left((I_d-CC^\top)BDB^\top
                           (I_d-CC^\top)\right).
\]
If $D\succ0$, the global minimizers therefore satisfy
$\operatorname{span}(C)=\operatorname{span}(B)$. Thus population ERM identifies
the desired representation. With finite samples, one can instead train the
factorized empirical objective by a first-order method, or minimize over each
$b_t$ in closed form and optimize the resulting profile objective over
$C$. Establishing that either procedure reaches a statistically accurate
solution requires an optimization-landscape analysis in addition to a
generalization argument; related results are described in the bibliographical
notes.

Instead, the linear-Gaussian model supplies a more specific route. The
population second moment of $Y^2XX^\top$ has the shared subspace as its leading
$k$-dimensional eigenspace. We estimate that moment and compute its leading
eigenvectors. The spectral estimator is therefore a deliberate choice: it is
an explicit polynomial-time procedure whose source-data dependence and
subspace error can be analyzed directly. The following calculations show both
why the source tasks identify the joint representation and how the statistical
saving survives after fitting the target head.

\subsubsection{Source diversity controls target representation error}

We first ask which source tasks constrain the representation needed by the
target. The Gaussian input model makes representation error a projection
error, so the comparison reduces to a relation between the task coefficients.

For $u=Cb-Ba_t$, the prediction residual is $X^\top u-\varepsilon$. The model
assumptions imply
\[
 \E[XX^\top]=I_d,
 \qquad
 \E[\varepsilon X]=\E[\varepsilon]\E[X]=0,
\]
where the second identity uses the independence of $\varepsilon$ and $X$.
Consequently,
\[
 R_t(g_b\circ h_C)
 =\E[(X^\top u-\varepsilon)^2]
 =u^\top\E[XX^\top]u-2u^\top\E[\varepsilon X]+\E[\varepsilon^2]
 =\|Cb-Ba_t\|_2^2+\E[\varepsilon^2].
\]
Since $C=B$ is allowed, $R_t^\star=\E[\varepsilon^2]$.
Minimizing over $b$ gives
\begin{equation}
 A_t(C)=\|(I_d-CC^\top)Ba_t\|_2^2.
 \label{eq:linear-representation-error}
\end{equation}
Define
\begin{equation}
 D=\frac1T\sum_{t=1}^T a_ta_t^\top,
 \qquad E_C=(I_d-CC^\top)B.
 \label{eq:task-diversity-matrix}
\end{equation}
Then $\overline A(C)=\operatorname{tr}(E_CDE_C^\top)$.

\begin{proposition}[Linear source-to-target comparison]
\label{prop:linear-calibration}
If $D\succ0$, every candidate $C$ satisfies
\begin{equation}
 A_0(C)\le\kappa_0\,\overline A(C),
 \qquad
 \kappa_0=a_0^\top D^{-1}a_0
 \le\frac{\|a_0\|_2^2}{\lambda_{\min}(D)}.
 \label{eq:linear-transfer-factor}
\end{equation}
\end{proposition}

\begin{proof}
Cauchy--Schwarz applied to $D^{-1/2}a_0$ and $D^{1/2}v$ gives
$(v^\top a_0)^2\le(a_0^\top D^{-1}a_0)v^\top Dv$ for every $v$.
Apply this inequality to each row of $E_C$ and add:
$\|E_Ca_0\|_2^2\le\kappa_0\operatorname{tr}(E_CDE_C^\top)$.
Finally, $D\succeq\lambda_{\min}(D)I_k$ implies
$D^{-1}\preceq\lambda_{\min}(D)^{-1}I_k$.
\end{proof}

The factor depends on the target direction as well as the source tasks.
A target aligned with a well-covered direction can have a small factor even
when other directions are poorly covered.\endnote{Positive definiteness of
$D$ is sufficient for all targets. For a fixed target with
$a_0\in\operatorname{range}(D)$, the same proof on that range gives
$\kappa_0=a_0^\top D^\dagger a_0$, where $D^\dagger$ is the Moore--Penrose
inverse. Coverage of directions unused by the target is unnecessary.}
If $\|a_0\|_2\le L$, the bound
$\Psi(u)\le L^2u/\lambda_{\min}(D)$ holds simultaneously for all such
targets.
The matrix $D$ gives a concrete measure of how strongly the source tasks
constrain the shared representation.

\subsubsection{Estimating the representation from source data}

We estimate the shared subspace using a second moment of the source
observations. The resulting \emph{spectral estimator} forms a representation
from leading eigenvectors of a data-dependent matrix. For the sample bound
below, additionally assume
$\varepsilon_{tj}\sim\cN(0,\sigma^2)$ on the source tasks. The moment
identity itself only uses centered noise of variance $\sigma^2$ independent
of the covariates.

The estimator belongs to the spectral family of representation-learning
methods. In linear models, the representation is a subspace and the estimator
recovers that subspace as the leading eigenspace of a covariance or moment
matrix. Kernel PCA, Laplacian eigenmaps, and spectral clustering replace the
linear moment matrix by a kernel matrix, graph Laplacian, or population
operator. Their coordinates can depend nonlinearly on the original input,
while the learned representation remains an eigenspace or a collection of
eigenfunctions. Spectral methods therefore cover a broad low-rank and
operator-based family. The bibliographical notes place this family among
other representation-learning objectives.

\begin{proposition}[Gaussian task moment]
\label{prop:gaussian-task-moment}
For one observation from a task with parameter $Ba$,
\begin{equation}
 \E[Y^2XX^\top]
 =(\|a\|_2^2+\sigma^2)I_d+2Baa^\top B^\top.
 \label{eq:gaussian-moment}
\end{equation}
Consequently the average population moment is
\begin{equation}
 M_\star=(\operatorname{tr}D+\sigma^2)I_d+2BDB^\top.
 \label{eq:population-moment}
\end{equation}
If $D\succ0$, its leading $k$-dimensional eigenspace is
$\operatorname{span}(B)$, with eigengap $2\lambda_{\min}(D)$.
\end{proposition}

\begin{proof}
Put $u=Ba$. Since $\varepsilon$ is independent of $X$ and has mean zero,
\[
 \E[2\varepsilon(X^\top u)XX^\top]
 =2\E[\varepsilon]\E[(X^\top u)XX^\top]=0,
 \qquad
 \E[\varepsilon^2XX^\top]
 =\E[\varepsilon^2]\E[XX^\top]=\sigma^2I_d.
\]
For standard Gaussian coordinates,
\[
 \E[X_iX_jX_rX_s]
 =\mathbf1_{\{i=j\}}\mathbf1_{\{r=s\}}
  +\mathbf1_{\{i=r\}}\mathbf1_{\{j=s\}}
  +\mathbf1_{\{i=s\}}\mathbf1_{\{j=r\}}.
\]
Expanding $(X^\top u)^2$ therefore gives
$\E[(X^\top u)^2XX^\top]=\|u\|_2^2I_d+2uu^\top$.
Since $u=Ba$ and $B^\top B=I_k$,
\[
 \|u\|_2^2=a^\top B^\top Ba=\|a\|_2^2,
 \qquad
 uu^\top=Baa^\top B^\top.
\]
Combining the three terms in $Y^2=(X^\top u+\varepsilon)^2$ now proves
\cref{eq:gaussian-moment}.
Average this identity over the source tasks:
\begin{align*}
 \frac1T\sum_{t=1}^T\E[Y_{t1}^2X_{t1}X_{t1}^\top]
 &=
 \left(\frac1T\sum_{t=1}^T\|a_t\|_2^2+\sigma^2\right)I_d
 +2B\left(\frac1T\sum_{t=1}^Ta_ta_t^\top\right)B^\top\\
 &=(\operatorname{tr}D+\sigma^2)I_d+2BDB^\top,
\end{align*}
because $T^{-1}\sum_t\|a_t\|_2^2=\operatorname{tr}D$. This proves
\cref{eq:population-moment}. Put $c=\operatorname{tr}D+\sigma^2$.
If $Dq=\lambda_i(D)q$, then
\[
 M_\star(Bq)=cBq+2BDq=(c+2\lambda_i(D))Bq.
\]
For $v\perp\operatorname{span}(B)$, the identity $B^\top v=0$ gives
$M_\star v=cv$. Thus the $k$ eigenvalues on $\operatorname{span}(B)$ are
$c+2\lambda_i(D)$, every eigenvalue on its orthogonal complement is $c$,
and the eigengap is $2\lambda_{\min}(D)$.
\end{proof}

The estimator forms
\begin{equation}
 \widehat M=\frac1N\sum_{t=1}^T\sum_{j=1}^n
                 Y_{tj}^2X_{tj}X_{tj}^\top
 \label{eq:empirical-moment}
\end{equation}
and takes its leading $k$ orthonormal eigenvectors as the columns of
$\widehat B$. Computing the moment and a dense eigendecomposition costs
$O(Nd^2+d^3)$ arithmetic operations. The procedure consists of one matrix
accumulation followed by one eigendecomposition.

Here is an elementary finite-sample guarantee that makes the statistical and
computational claim explicit. Sharper analyses improve its dimension
dependence.

\paragraph{Davis--Kahan bound used below.}
Let $M$ and $\widehat M=M+F$ be symmetric matrices. Suppose the leading
$k$-dimensional eigenspace of $M$ has orthonormal basis $B$ and eigengap
\[
 \gamma=\lambda_k(M)-\lambda_{k+1}(M)>0.
\]
If $\widehat B$ contains leading $k$ orthonormal eigenvectors of $\widehat M$,
then the Davis--Kahan theorem gives
\begin{equation}
 \|(I_d-\widehat B\widehat B^\top)B\|_{\rm op}
 \le
 \min\left\{1,\frac{2\|F\|_{\rm op}}{\gamma}\right\}.
 \label{eq:davis-kahan-bound}
\end{equation}
The left side is the sine of the largest principal angle between the two
eigenspaces.

\begin{theorem}[A finite-sample spectral guarantee]
\label{thm:spectral-guarantee}
Suppose $\lambda=\lambda_{\min}(D)>0$, and let
\[
 V=\max_{t\in[T]}\E[Y_{t1}^4\|X_{t1}\|_2^4].
\]
The spectral estimator satisfies
\begin{equation}
 \E[A_0(\widehat B)]
 \le\frac{\|a_0\|_2^2 V}{N\lambda^2}.
 \label{eq:spectral-target-bound}
\end{equation}
In the Gaussian model, if
$s^2=\max_t(\|a_t\|_2^2+\sigma^2)$, then
\[
 V\le\sqrt{105}\,s^4\sqrt{d(d+2)(d+4)(d+6)}.
\]
\end{theorem}

\begin{proof}
The goal is to control the target representation error
$\E[A_0(\widehat B)]$. The representation-error identity
\cref{eq:linear-representation-error}, which also starts the proof of
\cref{prop:linear-calibration}, gives
\begin{align}
 A_0(\widehat B)
 &=\|(I_d-\widehat B\widehat B^\top)Ba_0\|_2^2 \notag\\
 &\le
 \|a_0\|_2^2
 \|(I_d-\widehat B\widehat B^\top)B\|_{\rm op}^2.
 \label{eq:target-error-from-subspace}
\end{align}
Thus it suffices to control the angle between
$\operatorname{span}(\widehat B)$ and $\operatorname{span}(B)$.

Proposition~\ref{prop:gaussian-task-moment} identifies
$\operatorname{span}(B)$ as the leading eigenspace of $M_\star$ and gives
the population eigengap
\[
 \gamma=2\lambda,\qquad \lambda=\lambda_{\min}(D).
\]
Put $F=\widehat M-M_\star$. Applying
\cref{eq:davis-kahan-bound} with this value of $\gamma$ gives
\[
 \|(I_d-\widehat B\widehat B^\top)B\|_{\rm op}
 \le
 \min\left\{1,\frac{\|F\|_{\rm op}}{\lambda}\right\}
 \le\frac{\|F\|_{\rm F}}{\lambda}.
\]
Substitution into \cref{eq:target-error-from-subspace} yields
\begin{equation}
 A_0(\widehat B)
 \le\frac{\|a_0\|_2^2}{\lambda^2}\|F\|_{\rm F}^2.
 \label{eq:target-error-from-moment}
\end{equation}
The remaining task is therefore to prove
$\E[\|F\|_{\rm F}^2]\le V/N$.

For this purpose, define
\[
 W_{tj}=Y_{tj}^2X_{tj}X_{tj}^\top,
 \qquad
 Q_{tj}
 =W_{tj}-\E[W_{tj}].
\]
Proposition~\ref{prop:gaussian-task-moment} and $N=nT$ give
\[
 \E[\widehat M]
 =\frac1N\sum_{t=1}^T\sum_{j=1}^n\E[W_{tj}]
 =\frac1T\sum_{t=1}^T\E[W_{t1}]
 =M_\star.
\]
Consequently, $F=N^{-1}\sum_{t,j}Q_{tj}$.
The matrices $Q_{tj}$ are independent and satisfy $\E[Q_{tj}]=0$. Hence,
for $(t,j)\ne(s,r)$,
\[
 \E[\langle Q_{tj},Q_{sr}\rangle_{\rm F}]
 =\langle\E[Q_{tj}],\E[Q_{sr}]\rangle_{\rm F}=0.
\]
Expanding the squared Frobenius norm therefore gives
\begin{align*}
 \E[\|F\|_{\rm F}^2]
 &=\frac1{N^2}\sum_{t,j}
   \E[\|Q_{tj}\|_{\rm F}^2]
 &\le\frac1{N^2}\sum_{t,j}\E[\|W_{tj}\|_{\rm F}^2]\\
 &=\frac1{N^2}\sum_{t,j}
   \E[Y_{tj}^4\|X_{tj}X_{tj}^\top\|_{\rm F}^2]\\
 &=\frac1{N^2}\sum_{t,j}
   \E[Y_{tj}^4\|X_{tj}\|_2^4]
 \le\frac VN.
\end{align*}
The inequality uses
$\E[\|W-\E[W]\|_{\rm F}^2]\le\E[\|W\|_{\rm F}^2]$, and the penultimate
identity uses
$\|xx^\top\|_{\rm F}^2=\operatorname{tr}(xx^\top xx^\top)=\|x\|_2^4$.
Taking expectations in \cref{eq:target-error-from-moment} now proves
\cref{eq:spectral-target-bound}.

It remains to bound $V$ under Gaussian noise. For each source task,
$Y_{t1}=X_{t1}^\top Ba_t+\varepsilon_{t1}$ is centered Gaussian with variance
$\|a_t\|_2^2+\sigma^2\le s^2$. Hence
$\E[Y_{t1}^8]\le105s^8$. Cauchy--Schwarz gives
$\E[Y_{t1}^4\|X_{t1}\|_2^4]
\le\sqrt{\E[Y_{t1}^8]\E[\|X_{t1}\|_2^8]}$.
Since $\E[\|X_{t1}\|_2^8]=d(d+2)(d+4)(d+6)$, we obtain
\[
 \E[Y_{t1}^4\|X_{t1}\|_2^4]
 \le
 \sqrt{105}\,s^4\sqrt{d(d+2)(d+4)(d+6)}.
\]
Taking the maximum over $t$ proves the stated bound on $V$.
\end{proof}

\subsubsection{Fitting the target head and bounding prediction risk}

With $\widehat B$ fixed, target training is ordinary least squares in $k$
features. Let $Z_0\in\R^{m\times k}$ have rows
$(\widehat B^\top X_{0j})^\top$ and let
$\boldsymbol Y_0=(Y_{01},\ldots,Y_{0m})^\top$. For $m>k+1$, set
\begin{equation}
 \widehat a_0=(Z_0^\top Z_0)^{-1}Z_0^\top\boldsymbol Y_0,
 \qquad \widehat\theta=\widehat B\widehat a_0.
 \label{eq:linear-target-head}
\end{equation}
Computing the target features and fitting the head costs
$O(mdk+mk^2+k^3)$ arithmetic operations. The exact Gaussian calculation in
\cref{app:gaussian-transfer} converts the representation bound into a
guarantee for the final target predictor.

\begin{corollary}[Linear transfer guarantee]
\label{cor:linear-transfer}
Under the assumptions of \cref{thm:spectral-guarantee}, let
$\widehat B$ be the spectral estimate and fit the target head by
\cref{eq:linear-target-head} using $m>k+1$ target examples.
With $\lambda$ and $V$ as in that theorem, the expected target excess risk is
bounded by
\begin{equation}
 \E[R_0(\widehat f_0)-R_0^\star]
 \le
 \left(1+\frac{k}{m-k-1}\right)
 \frac{\|a_0\|_2^2V}{N\lambda^2}
 +\sigma_0^2\frac{k}{m-k-1}.
 \label{eq:computational-transfer}
\end{equation}
\end{corollary}

\begin{proof}
The target sample is independent of $\widehat B$. Conditional on
$\widehat B$, \cref{thm:exact-transfer-risk} gives
\[
 \E[\|\widehat\theta-Ba_0\|_2^2\mid\widehat B]
 =\left(1+\frac{k}{m-k-1}\right)A_0(\widehat B)
  +\sigma_0^2\frac{k}{m-k-1}.
\]
Average over the source data, apply \cref{thm:spectral-guarantee}, and use
\cref{eq:linear-target-objective}.
\end{proof}

The first term accounts for learning the representation from source data;
the second accounts for fitting the head from target data. For comparison,
ordinary least squares using only the target sample in $\R^d$ has expected
excess risk $\sigma_0^2d/(m-d-1)$ when $m>d+1$. Accurate source-based
features reduce the dimension in the target estimation term from $d$ to $k$.
Together with the arithmetic costs above, \cref{cor:linear-transfer} gives a
computational and statistical guarantee for transfer.

\paragraph{Connection to neural representations.}
A downstream task can use a neural network's frozen features and a newly
fitted head. The transfer analyses explain what a guarantee for that use
requires: source training must control the target representation error, and
target data must suffice to fit the head. For language-model pretraining,
the remaining question is how next-token loss constrains downstream
representation error.

\clearpage
\section{Take-home messages}

\begin{itemize}[leftmargin=*]
\item Joint ERM pays the representation complexity once across the source
tasks and the head complexity once per task. The shared boxes example gives
a strict statistical advantage over independent learning.
\item Misspecification contributes the sharing cost $\Delta_{\rm share}$.
Reduced estimation error must compensate for any added approximation error.
\item A bound on $A_0(h)$ in terms of $\overline A(h)$ converts source
learning into a guarantee for a particular target. The examples verify the
relation and exhibit a case where it fails.
\item With independent target examples and a frozen representation, head
fitting pays only the head-class complexity. Simultaneous individual-task
guarantees require suitable source-to-task relations for those tasks.
\item In the Gaussian linear model, a spectral representation estimate
followed by target least squares has an explicit computation cost and a
finite-sample target-risk guarantee.
\end{itemize}

\printendnotes

\section*{Bibliographical notes}

For a broad introduction, \citet{bengio2013} survey representation learning,
and Chapter 15 of \citet{goodfellow2016} gives a textbook treatment.
The following map organizes several major families by the information used to
learn the representation and by the property retained in the representation.

\begin{description}[style=unboxed,leftmargin=1.7em,labelindent=0pt]
\item[Linear, latent, and spectral structure.]
PCA retains directions of large variance; factor models and ICA seek latent
coordinates with additional statistical structure; graph and kernel methods
use eigenvectors or eigenfunctions of similarity operators. The review by
\citet{jolliffe2016}, the ICA book by \citet{hyvarinen2001}, and the spectral
clustering tutorial by \citet{vonluxburg2007} provide readable introductions.
\citet{belkin2003} show how a graph Laplacian produces nonlinear coordinates
for data near a manifold.

\item[Supervised task sharing and transfer.]
Labels from several tasks train features that support task-specific heads.
Multitask learning, learning-to-learn, and the linear model in this lecture
belong to this family. The statistical sources most closely connected to the
lecture are discussed below.

\item[Reconstruction and generative modeling.]
Autoencoders, probabilistic latent-variable models, and related generative
methods learn coordinates that retain information needed to reconstruct or
model the observations. \citet{bengio2013} and \citet{goodfellow2016} connect
these objectives to deep representation learning.

\item[Predictive and contrastive self-supervision.]
Context prediction, masked or next-token prediction, and contrastive
objectives derive training signals from the observations themselves.
\citet{lekhac2020} give an organized review of the contrastive branch.

\item[Representations for sequential decisions.]
Control and reinforcement-learning methods seek state variables that preserve
reward-relevant and transition-relevant information. The overview by
\citet{lesort2018} organizes reconstruction, forward-model, inverse-model, and
robotic-prior approaches.

\item[Geometric and causal structure.]
Equivariant methods encode symmetries and relational structure; causal
representation learning seeks variables that support interventions and
distributional changes. \citet{bronstein2021} survey geometric deep learning,
and \citet{scholkopf2021} survey causal representation learning.
\end{description}

Multitask feature learning and transfer have a substantial statistical
literature for general classes of representations.
\citet{baxter2000} studies learning a hypothesis space from tasks sampled
from an environment. \citet{maurer2016} gives general representation-class
bounds for both observed-task prediction and learning new tasks, with separate
costs for representations and task-specific predictors.
The finite-class results here are elementary versions of this complexity
separation, proved using the tools already developed in the course.

The closest source for bounds such as \cref{eq:linear-calibration} is the
task-diversity condition of \citet{tripuraneni2020}, Definition 3. Their
condition bounds a worst-case target representation discrepancy by an average
source discrepancy divided by a diversity parameter, plus an additive error.
Our fixed-target formulation uses the transfer envelope $\Psi$ and
class-relative approximation errors; the displayed finite-class theorem is
proved directly. The transfer envelope $\Psi$ expresses the connection to the
fixed source tasks directly.

The Gaussian linear model and moment estimator follow
\citet{tripuraneni2021}. The simple Frobenius-moment bound in
\cref{thm:spectral-guarantee} is deliberately weaker than sharp analyses of
this model and provides a complete finite-sample argument. The eigenspace
perturbation step is the Davis--Kahan theorem \citep{daviskahan1970}.
The inverse-Gram expectation in \cref{eq:inverse-wishart-mean} is the
standard inverse-Wishart mean identity.

\citet{tripuraneni2021} also study optimization-based representation
learning. They use a Burer--Monteiro factorization of the task-parameter
matrix and add a balancing regularizer. Under their sub-Gaussian design,
task-diversity, incoherence, sample-size, and uniform-task-sampling
conditions, every local minimum in a specified interior region recovers the
shared subspace up to statistical error. Their local-minimum result, together
with general results on escaping strict saddles, gives a route from first-order
optimization to an accurate representation. Their theorem applies to the
regularized factorization in its stated constraint region.

Task-boundary terminology is illustrated by \citet{aljundi2019}.
The optional intrinsic-dimension comparison in
\cref{app:intrinsic-dimension} uses \citet{kohler2026},
\citet{kpotufe2011}, \citet{kpotufe2013}, and \citet{farahmand2007}.

\section*{Glossary}

\begin{description}[style=unboxed,leftmargin=1.5em,labelindent=0pt]
\item[Risk notation]
A sample in the subscript marks an empirical risk, so $R_t$ uses the law
$P_t$ and $R_{S_t}$ uses the sample $S_t$. A bar averages over the $T$ source
tasks, giving $\overline R$ and $\overline R_S$. A hat marks a quantity
returned by a learner, such as $\widehat h$ or $\widehat B$. The subscripts
$\mathrm{src}$ and $\mathrm{trg}$ distinguish source- and target-side
quantities.
\item[Representation-learning algorithm]
A map $\cA:(\cX\times\cY)^n\to\cH\times\cG$ returning a representation
and a head for supervised prediction.
\item[Representation and head]
Maps $h:\cX\to\cZ$ and $g:\cZ\to\cY$, whose composition predicts a
label. In the multitask setting, $h$ is shared and each task has its own head.
\item[Multitask learning]
Jointly learning predictors for several tasks and evaluating on those tasks.
\item[Transfer learning]
Using information learned from source tasks to learn a predictor for a target
task; here, fitting a new head on a frozen source-trained representation.
\item[Learning-to-learn / meta-learning]
Learning how to fit new tasks from experience on source tasks.
\item[Joint ERM]
Minimization of average empirical task risk over one shared representation
and one head for each task.
\item[Representation error $A_t(h)$]
The extra population risk imposed on task $t$ by using $h$, after optimizing
the head, relative to the best pair in $\cH\times\cG$.
\item[Sharing cost $\Delta_{\rm share}$]
The minimum average representation error over one shared $h$.
\item[Transfer envelope $\Psi$]
The largest target representation error among candidate representations
whose average source representation error is at most a specified value.
\item[Task-diversity matrix]
The matrix $D=T^{-1}\sum_t a_ta_t^\top$ in the linear model.
Its quadratic form measures source coverage of a coefficient direction.
\item[Spectral representation estimator]
An estimator that uses leading eigenvectors of a data-dependent matrix to
define the feature map.
\item[Negative transfer]
The excess risk of a sharing or transfer procedure relative to a specified
nonsharing baseline.
\item[Intrinsic dimension]
A dimension describing the geometry of an input distribution, rather than
the complexity of a shared family of tasks.
\end{description}

\clearpage
\section{Exercises}
\label{sec:exercises}

\exerciseinstructions
\setcounter{courseexercise}{0}
\exercisegroup{Joint learning and transfer}
\begin{exercisechunk}
\item \textbf{Finite-class zero-loss bound.}
\label[exercise]{ex:l13-bounded-loss-consistency}
\exerciseprofile{2}{2}{2}{\exproof\quad\excalculation\quad\exinterpretation}{%
Extend the finite-class consistency argument from classification to bounded
losses and independent observations with different laws.}
Use the setup immediately preceding
\cref{prop:independent-finite-class-consistency}.
\begin{enumerate}[(a),beginpenalty=10000]
\item Let $Z$ be a random variable with values in $[0,1]$. Prove
\[
 \P(Z=0)\le 1-\E[Z].
\]
Show that equality holds when $Z$ takes values in $\{0,1\}$.
\item For a fixed $f\in\cF$, use independence and part (a) to prove
\[
 \P\bigl(\ell(f(U_i),V_i)=0\text{ for every }i\in[m]\bigr)
 \le e^{-m\overline R_m(f)}.
\]
\item Use a union bound to prove the probability statement in
\cref{prop:independent-finite-class-consistency}. Your argument must apply to
every choice among zero-training-loss members of $\cF$, including a randomized
choice.
\item Integrate the resulting tail bound to prove the expected-risk statement.
Explain where the assumption that some $f^\star\in\cF$ has
$R_i(f^\star)=0$ for every $i$ is used.
\end{enumerate}
\end{exercisechunk}

\begin{exercisechunk}
\item \textbf{The approximation cost of sharing.}
\label[exercise]{ex:l13-sharing-cost}
\exerciseprofile{2}{2}{2}{\excalculation\quad\exproof\quad\exinterpretation}{%
Compute a sharing cost when every task is individually realizable, and
compare its effect with the estimation cost.}
Let $U,V$ be independent fair bits and let $X=(U,V)$ on every task.
A fraction $p$ of the $T$ source tasks have label $U$, and the remaining
fraction have label $V$, where $pT$ is an integer.
Take $\cH=\{h_U,h_V\}$ with $h_U(U,V)=U$ and $h_V(U,V)=V$,
$\cG=\{z\mapsto z,\ z\mapsto1-z\}$, and $0$--$1$ loss.
Each task supplies $n$ independent examples, and the task samples are
independent.
\begin{enumerate}[(a),beginpenalty=10000]
\item Compute $R_t^\star$, $A_t(h_U)$, and $A_t(h_V)$ for both kinds of
task. Deduce the exact value of $\Delta_{\rm share}$.
\item Specialize the agnostic joint-ERM bound to these classes, assuming
exact empirical optimization.
\item Suppose $0<p<1$. Explain why increasing $n$ cannot make the average
risk of a shared predictor tend to zero. Compare with consistent independent
learning. Identify the assumption of the realizable joint zero-loss
theorem that fails.
\end{enumerate}
\end{exercisechunk}

\begin{exercisechunk}
\item \textbf{Sharpness of the transfer envelope.}
\label[exercise]{ex:l13-transfer-envelope}
\exerciseprofile{2}{1}{2}{\exproof\quad\exinterpretation}{%
Show that the transfer envelope is the strongest target guarantee implied
only by an average source representation-error bound.}
Let $\Psi$ be the transfer envelope from \cref{eq:transfer-envelope}.
\begin{enumerate}[(a),beginpenalty=10000]
\item Prove that $\Psi$ is nondecreasing and that
$A_0(h)\le\Psi(\overline A(h))$ for every $h\in\cH$.
\item Let $\phi:[0,\infty)\to[0,\infty)$ be nondecreasing and satisfy
$A_0(h)\le\phi(\overline A(h))$ for every $h\in\cH$.
Prove that $\Psi(u)\le\phi(u)$ for every $u\ge0$.
When $\{h:\overline A(h)\le u\}$ is nonempty, prove that some feasible
$h_u$ attains $A_0(h_u)=\Psi(u)$.
Explain the precise worst-case statement established by these facts.
\item In the target-feature example of \cref{sec:calibration-examples},
compute $\Psi(0)$ and identify a maximizing representation.
Give a joint source-ERM tie-breaking rule that always selects this
representation. Explain why the value of $\Psi(0)$ does not imply that every
source learner selects it.
\end{enumerate}
\end{exercisechunk}

\begin{exercisechunk}
\item \textbf{Coverage of a target and simultaneous task guarantees.}
\label[exercise]{ex:l13-transfer-coverage}
\exerciseprofile{2}{2}{2}{\exproof\quad\exinterpretation}{%
Turn a quantitative source-to-target relation into useful guarantees for
designated individual tasks.}
Use the finite-class bounded-loss setup of the lecture, and suppose
$\P\{\overline A(\widehat h)\le\varepsilon_{\rm src}\}
\ge1-\delta_{\rm src}$.
\begin{enumerate}[(a),beginpenalty=10000]
\item Suppose at least $\rho T$ source tasks have the same law as a fixed
target $P_0$, where $0<\rho\le1$. Prove that
$A_0(h)\le\overline A(h)/\rho$ for every $h$.
Explain why the factor depends on the fraction of relevant source tasks,
rather than their number alone.
\item More generally, for targets $j=1,\ldots,q$, suppose known constants
$\kappa_j,\beta_j\ge0$ satisfy
$A_j(h)\le\kappa_j\overline A(h)+\beta_j$ for all $h$.
Each target supplies $m_j$ fresh i.i.d.\ examples independent of the source
data. Freeze $\widehat h$ and fit each head by exact ERM.
Give a simultaneous bound on all $q$ excess risks with probability at least
$1-\delta_{\rm src}-\delta_{\rm heads}$, assigning each head the same failure
probability. Are the target datasets required to be independent of one
another for your proof?
\end{enumerate}
\end{exercisechunk}

\begin{exercisechunk}
\item \textbf{A small target sample can select the head.}
\label[exercise]{ex:l13-label-reversal}
\exerciseprofile{3}{2}{3}{\exproof\quad\excalculation\quad\exinterpretation}{%
Sharpen the generic target-ERM bound when only the orientation of a shared
Boolean feature remains to be learned.}
Use the label-reversal model of \cref{eq:label-reversal-model}.
Suppose the source procedure has an event of probability at least
$1-\delta_{\rm src}$ on which
$r:=A_0(\widehat h)\le\varepsilon\le1/4$.
Freeze $\widehat h$ and select one of the two heads by ERM on $m$
fresh target examples independent of all source data.
\begin{enumerate}[(a),beginpenalty=10000]
\item Conditional on a source sample in that event, prove that the
probability of choosing the higher-risk head is at most
$\exp(-m/8)$. Allow arbitrary tie breaking.
\item Deduce a guarantee that the final target risk is at most $\varepsilon$
with probability at least $1-\delta_{\rm src}-\delta_{\rm trg}$.
Give sufficient $n$ and $m$ when the source learner is joint ERM with
$n$ examples from each of $T$ tasks.
\item For fixed $\varepsilon\in(0,1/4)$ and fixed confidence, compare the
target sample requirement with the boxes lower bound for a learner using
only target data. Explain which information the source data supply and which
information the target data supply.
\end{enumerate}
\begin{hints}
For each target example, the two heads' mistake indicators sum to one.
Apply Hoeffding's inequality to the empirical risk of the better head.
\end{hints}
\end{exercisechunk}

\exercisegroup{Linear representations}
\begin{exercisechunk}
\item \textbf{Many repeated tasks can hide a shared feature.}
\label[exercise]{ex:l13-task-diversity}
\exerciseprofile{2}{1}{2}{\exproof\quad\excalculation\quad\exinterpretation}{%
Distinguish the number of source tasks from coverage of the shared feature
directions.}
Let $d=3$, $k=2$, and let every source-task head have the form
$a_i=(c_i,0)^\top$, where $c_i\in\R$. Consider the two feature matrices
\[
 B=[e_1\ e_2]
 \qquad\text{and}\qquad
 \widetilde B=[e_1\ e_3].
\]
For either feature matrix, source task $i$ generates
$Y=X^\top Ba_i+\varepsilon$, where $X\sim\cN(0,I_3)$ and the same independent
noise law is used in both models.
\begin{enumerate}[(a),beginpenalty=10000]
\item Prove that the joint law of $(X,Y)$ for every source task is identical
under $B$ and $\widetilde B$.
\item For uniform sampling from $M$ source tasks, compute the task-diversity
matrix
\[
 D=\frac1M\sum_{i=1}^M a_ia_i^\top
\]
and its eigenvalues.
\item Let the new-task head be $a_{\rm new}=(0,1)^\top$. Compute its ambient
parameter under each feature matrix. Explain why no procedure using only the
source data can determine which of these two new-task parameters is correct.
\end{enumerate}
\end{exercisechunk}

\begin{exercisechunk}
\item \textbf{The angle between tasks controls the eigengap.}
\label[exercise]{ex:l13-moment-eigenspace}
\exerciseprofile{2}{2}{2}{\excalculation\quad\exinterpretation}{%
Relate geometric diversity among two task heads to stability of population
moment recovery.}
Let $k=2$, and sample two source tasks uniformly. Their unit-norm heads are
\[
 a_1=\begin{pmatrix}1\\0\end{pmatrix},
 \qquad
 a_2=\begin{pmatrix}\cos\alpha\\\sin\alpha\end{pmatrix},
 \qquad 0<\alpha\le\frac\pi2.
\]
\begin{enumerate}[(a),beginpenalty=10000]
\item Compute
$D=(a_1a_1^\top+a_2a_2^\top)/2$ and show that its eigenvalues are
\[
 \lambda_{\max}(D)=\frac{1+\cos\alpha}{2},
 \qquad
 \lambda_{\min}(D)=\frac{1-\cos\alpha}{2}.
\]
\item Under the Gaussian model of
\cref{prop:gaussian-task-moment}, show that the gap between the second and
third eigenvalues of $M_\star$ is $1-\cos\alpha$.
\item Determine the leading-order behavior of this eigengap as
$\alpha\downarrow0$. Explain why collecting more observations from two nearly
parallel tasks is statistically different from observing two well-separated
tasks.
\end{enumerate}
\begin{hints}
The trace and determinant of $D$ determine its two eigenvalues. Use
$1-\cos\alpha\sim\alpha^2/2$ for the last part.
\end{hints}
\end{exercisechunk}

\begin{exercisechunk}
\item \textbf{Projection error is representation bias.}
\label[exercise]{ex:l13-transfer-projection}
\exerciseprofile{2}{2}{2}{\exproof\quad\excalculation\quad\exinterpretation}{%
Derive the deterministic approximation term that enters the new-task transfer
risk.}
Let $B,\widehat B\in\R^{d\times k}$ have orthonormal columns, and define
$P_B=BB^\top$, $\widehat P=\widehat B\widehat B^\top$, and
\[
 \delta=\|(I_d-\widehat P)B\|_{\mathrm{op}}.
\]
Fix an ambient task parameter $\theta\in\R^d$.
\begin{enumerate}[(a),beginpenalty=10000]
\item Show that
\[
 \min_{b\in\R^k}\|\widehat Bb-\theta\|_2^2
 =\|(I_d-\widehat P)\theta\|_2^2,
\]
and identify a minimizing vector $b$.
\item If $\theta=Ba$, prove that
\[
 \|(I_d-\widehat P)\theta\|_2^2
 \le\delta^2\|a\|_2^2.
\]
\item If $\theta=Ba+r$ with $B^\top r=0$, prove that
\[
 \|(I_d-\widehat P)\theta\|_2
 \le\delta\|a\|_2+\|r\|_2.
\]
Identify which term is caused by estimating the shared subspace and which
term is present even when $\widehat B=B$.
\end{enumerate}
\end{exercisechunk}

\begin{exercisechunk}
\item \textbf{An exact negative-transfer threshold.}
\label[exercise]{ex:l13-negative-transfer}
\exerciseprofile{2}{2}{2}{\exproof\quad\excalculation\quad\exinterpretation}{%
Determine when misspecification outweighs the variance saved by a known
low-dimensional representation.}
Assume the Gaussian new-task model of
\cref{thm:exact-transfer-risk}. Let $n>d+1$, let $k<d$, and suppose the
noise variance is $\sigma^2>0$ and the parameter is $\theta=Ba+r$,
where $B^\top B=I_k$, $B^\top r=0$, and the known
representation $B$ is used for transfer. Compare ordinary least squares in
$\operatorname{span}(B)$ with ordinary least squares in $\R^d$.
\begin{enumerate}[(a),beginpenalty=10000]
\item Use \cref{thm:exact-transfer-risk} to write the exact expected squared
parameter error of the representation-restricted estimator and of the ambient
estimator.
\item Prove that the representation-restricted estimator has smaller expected
error exactly when
\[
 \|r\|_2^2<\frac{\sigma^2(d-k)}{n-d-1}.
\]
\item Fix $r\ne0$. Explain how the comparison changes as $n$ grows. Also
explain what happens to ambient least squares when $k+1<n<d$ and when
$n\in\{d,d+1\}$. For the latter case, you may use the fact that the expected
trace of the inverse Gaussian Gram matrix is infinite.
\end{enumerate}
\end{exercisechunk}


\appendix
\renewcommand{\thesection}{13.\Alph{section}}
\section{Exact Gaussian head-estimation risk}
\label{app:gaussian-transfer}

This appendix gives the exact head-estimation term used in
\cref{eq:computational-transfer} and a concrete negative-transfer comparison.
Fix $\widehat B\in\R^{d\times k}$ with orthonormal columns and
$\theta\in\R^d$. Put
\begin{equation}
 \widehat P=\widehat B\widehat B^\top,
 \qquad
 A(\theta,\widehat B)=\|(I_d-\widehat P)\theta\|_2^2.
 \label{eq:projection-bias}
\end{equation}
Observe $n$ independent pairs with
\begin{equation}
 Y_j=X_j^\top\theta+\varepsilon_j,\qquad
 X_j\sim\cN(0,I_d),\qquad j\in[n].
 \label{eq:new-task-model}
\end{equation}
The noises are mutually independent and independent of all covariates,
have mean zero, and have common variance $\sigma^2$.
Let $Z$ have rows $(\widehat B^\top X_j)^\top$ and let
$Y=(Y_1,\ldots,Y_n)^\top$. Define
\begin{equation}
 \widehat a=(Z^\top Z)^{-1}Z^\top Y,
 \qquad \widehat\theta=\widehat B\widehat a.
 \label{eq:new-task-ols}
\end{equation}

\begin{theorem}[Exact Gaussian transfer risk]
\label{thm:exact-transfer-risk}
Under this model, if $n>k+1$, then
\begin{equation}
 \E[\|\widehat\theta-\theta\|_2^2]
 =A(\theta,\widehat B)
  +(\sigma^2+A(\theta,\widehat B))\frac{k}{n-k-1}.
 \label{eq:exact-transfer-risk}
\end{equation}
For an independent test covariate $X\sim\cN(0,I_d)$, the expected prediction
error is
\begin{equation}
 \E[(X^\top\widehat\theta-X^\top\theta)^2]
 =\E[\|\widehat\theta-\theta\|_2^2].
 \label{eq:prediction-risk-equality}
\end{equation}
\end{theorem}

\begin{proof}
Write $b_\star=\widehat B^\top\theta$ and
$v=(I_d-\widehat P)\theta$. Then
$\theta=\widehat Bb_\star+v$ with $\widehat B^\top v=0$, and
\[
 Y=Zb_\star+\xi,\qquad \xi_j=X_j^\top v+\varepsilon_j.
\]
For each $j$,
\[
 \operatorname{Cov}(\widehat B^\top X_j,X_j^\top v)
 =\widehat B^\top\E[X_jX_j^\top]v
 =\widehat B^\top v=0.
\]
The two projections are jointly Gaussian, so zero covariance makes them
independent. The noise $\varepsilon_j$ is also independent of $X_j$.
Therefore $\xi$ is independent of $Z$, has mean zero, and has covariance
$(\sigma^2+\|v\|_2^2)I_n$. Conditional on $Z$,
\[
 \E[\|\widehat a-b_\star\|_2^2\mid Z]
 =(\sigma^2+\|v\|_2^2)\operatorname{tr}((Z^\top Z)^{-1}).
\]
The Wishart inverse-mean identity gives, for $n>k+1$,
\begin{equation}
 \E[(Z^\top Z)^{-1}]=\frac{I_k}{n-k-1}.
 \label{eq:inverse-wishart-mean}
\end{equation}
Since $\widehat B(\widehat a-b_\star)$ is orthogonal to $v$,
\[
 \|\widehat\theta-\theta\|_2^2
 =\|\widehat a-b_\star\|_2^2+\|v\|_2^2.
\]
Take expectations. For the independent test covariate $X$, independence from
$\widehat\theta$ and
$\E[XX^\top]=I_d$ give
\[
 \E[(X^\top\widehat\theta-X^\top\theta)^2\mid\widehat\theta]
 =(\widehat\theta-\theta)^\top\E[XX^\top]
   (\widehat\theta-\theta)
 =\|\widehat\theta-\theta\|_2^2.
\]
\end{proof}

If $\widehat B$ is random but independent of the target sample, apply the
theorem conditionally on $\widehat B$ and then average. In the exact-sharing
model $\theta=Ba_0$, the resulting identity and
\cref{eq:spectral-target-bound} give \cref{eq:computational-transfer}.

For a known correct representation, $A=0$ and the risk is
$\sigma^2k/(n-k-1)$. Ambient ordinary least squares has risk
$\sigma^2d/(n-d-1)$ when $n>d+1$.
For misspecified sharing $\theta=Ba+r$ with $B^\top r=0$ and
$\widehat B=B$, the representation error is $\|r\|_2^2$.
Exercise~\ref{ex:l13-negative-transfer} asks you to derive the exact comparison:
the restricted estimator has smaller expected error if and only if
\[
 \|r\|_2^2<\frac{\sigma^2(d-k)}{n-d-1}.
\]
The restriction can save variance with little target data, but a fixed
nonzero residual eventually outweighs that saving.
\section{Optional comparison: intrinsic input dimension}
\label{app:intrinsic-dimension}

A single regression problem can become easier because of the input
distribution's geometry, without any sharing across tasks.
For a pair $(X,Y)$ with $\E[Y^2]<\infty$, let $P_X$ be the input law on
$\R^d$ and choose a measurable version $m$ of the regression function, so
$m(X)=\E[Y\mid X]$ almost surely.
If $P_X$ concentrates near a $d^\ast$-dimensional set, a regression rate may
depend on $d^\ast$ rather than $d$.

A theorem of Kohler, Molinero R\"omer, and Krzy\.{z}ak gives such a rate for a
specified deep sigmoid estimator: prescribed random initialization, random
pruning before training, output truncation, and full-batch gradient descent
with data-dependent Armijo-style stepsizes. Include the training sample and
the estimator's prescribed initialization and pruning in the same joint
experiment. Under their smoothness, tail,
architecture, and tuning assumptions, and support in a shrinking neighborhood
of a compact Lipschitz manifold, their bound has the form
\[
 \E[\|\widehat m_n-m\|_{L^2(P_X)}^2]
 \le C'(\log n)^{4d^\ast L+d+3}n^{-2p/(2p+d^\ast)}.
\]
Here $p$ is smoothness, $L$ is the prescribed depth parameter, and the
the tuning uses $p,d^\ast$. The tube radius is
$n^{-1/(2p+d^\ast)}(\log n)^{-4L-1}$, not fixed positive thickness.
The sufficient widths and iteration counts have large polynomial exponents.
This is a theorem for that procedure under those assumptions, not evidence
for the manifold hypothesis or a guarantee for ordinary network training.

Classical local methods also obtain intrinsic-dimension rates:
Kpotufe's nearest-neighbor result adapts to local dimension, and Kpotufe and
Garg's kernel result adapts to local dimension and local H\"older smoothness.
Farahmand, Szepesv\'ari, and Audibert developed a related manifold-adaptive
program. Their assumptions and smoothness ranges should be compared
explicitly, rather than treating all rates as one identical theorem.
The neural result's distinctive content is its guarantee for a specified
gradient-descent estimator. Low-dimensional input geometry and shared
representations across tasks are different statistical assumptions.

\clearpage
\begin{thebibliography}{99}
\raggedright

\bibitem[Aljundi et al.(2019)]{aljundi2019}
Rahaf Aljundi, Klaas Kelchtermans, and Tinne Tuytelaars.
\newblock Task-free continual learning.
\newblock In \emph{Proceedings of the IEEE/CVF Conference on Computer Vision
and Pattern Recognition}, pages 11254--11263, 2019.
\newblock \url{https://openaccess.thecvf.com/content_CVPR_2019/html/Aljundi_Task-Free_Continual_Learning_CVPR_2019_paper.html}.

\bibitem[Baxter(2000)]{baxter2000}
Jonathan Baxter.
\newblock A model of inductive bias learning.
\newblock \emph{Journal of Artificial Intelligence Research}, 12:149--198, 2000.
\newblock \url{https://doi.org/10.1613/jair.731}.

\bibitem[Belkin and Niyogi(2003)]{belkin2003}
Mikhail Belkin and Partha Niyogi.
\newblock Laplacian eigenmaps for dimensionality reduction and data
representation.
\newblock \emph{Neural Computation}, 15(6):1373--1396, 2003.
\newblock \url{https://doi.org/10.1162/089976603321780317}.

\bibitem[Bengio et al.(2013)]{bengio2013}
Yoshua Bengio, Aaron Courville, and Pascal Vincent.
\newblock Representation learning: A review and new perspectives.
\newblock \emph{IEEE Transactions on Pattern Analysis and Machine
Intelligence}, 35(8):1798--1828, 2013.
\newblock \url{https://doi.org/10.1109/TPAMI.2013.50}.

\bibitem[Bronstein et al.(2021)]{bronstein2021}
Michael M. Bronstein, Joan Bruna, Taco Cohen, and Petar Veli\v{c}kovi\'c.
\newblock Geometric deep learning: Grids, groups, graphs, geodesics, and
gauges.
\newblock 2021.
\newblock \url{https://arxiv.org/abs/2104.13478}.

\bibitem[Davis and Kahan(1970)]{daviskahan1970}
Chandler Davis and William M. Kahan.
\newblock The rotation of eigenvectors by a perturbation. III.
\newblock \emph{SIAM Journal on Numerical Analysis}, 7(1):1--46, 1970.
\newblock \url{https://doi.org/10.1137/0707001}.

\bibitem[Farahmand et al.(2007)]{farahmand2007}
Amir-massoud Farahmand, Csaba Szepesv\'ari, and Jean-Yves Audibert.
\newblock Manifold-adaptive dimension estimation.
\newblock In \emph{Proceedings of the 24th International Conference on
Machine Learning}, pages 265--272, 2007.
\newblock \url{https://doi.org/10.1145/1273496.1273530}.

\bibitem[Goodfellow et al.(2016)]{goodfellow2016}
Ian Goodfellow, Yoshua Bengio, and Aaron Courville.
\newblock \emph{Deep Learning}.
\newblock MIT Press, 2016.
\newblock \url{https://www.deeplearningbook.org/}.

\bibitem[Hyv\"arinen et al.(2001)]{hyvarinen2001}
Aapo Hyv\"arinen, Juha Karhunen, and Erkki Oja.
\newblock \emph{Independent Component Analysis}.
\newblock Wiley, 2001.
\newblock \url{https://doi.org/10.1002/0471221317}.

\bibitem[Jolliffe and Cadima(2016)]{jolliffe2016}
Ian T. Jolliffe and Jorge Cadima.
\newblock Principal component analysis: A review and recent developments.
\newblock \emph{Philosophical Transactions of the Royal Society A},
374(2065):20150202, 2016.
\newblock \url{https://doi.org/10.1098/rsta.2015.0202}.

\bibitem[Kohler et al.(2026)]{kohler2026}
Michael Kohler, Vincent Molinero R\"omer, and Adam Krzy\.{z}ak.
\newblock Learning of deep neural network regression estimates using gradient
descent with pruning.
\newblock 2026.
\newblock \url{https://arxiv.org/abs/2608.25743}.

\bibitem[Kpotufe(2011)]{kpotufe2011}
Samory Kpotufe.
\newblock $k$-NN regression adapts to local intrinsic dimension.
\newblock In \emph{Advances in Neural Information Processing Systems 24},
2011.
\newblock \url{https://proceedings.neurips.cc/paper/2011/hash/05f971b5ec196b8c65b75d2ef8267331-Abstract.html}.

\bibitem[Kpotufe and Garg(2013)]{kpotufe2013}
Samory Kpotufe and Vikas K. Garg.
\newblock Adaptivity to local smoothness and dimension in kernel regression.
\newblock In \emph{Advances in Neural Information Processing Systems 26},
2013.
\newblock \url{https://proceedings.neurips.cc/paper/2013/hash/28fc2782ea7ef51c1104ccf7b9bea13d-Abstract.html}.

\bibitem[Le-Khac et al.(2020)]{lekhac2020}
Phuc H. Le-Khac, Graham Healy, and Alan F. Smeaton.
\newblock Contrastive representation learning: A framework and review.
\newblock \emph{IEEE Access}, 8:193907--193934, 2020.
\newblock \url{https://doi.org/10.1109/ACCESS.2020.3031549}.

\bibitem[Lesort et al.(2018)]{lesort2018}
Timoth\'ee Lesort, Natalia D\'iaz-Rodr\'iguez, Jean-Fran\c{c}ois Goudou, and
David Filliat.
\newblock State representation learning for control: An overview.
\newblock \emph{Neural Networks}, 108:379--392, 2018.
\newblock \url{https://doi.org/10.1016/j.neunet.2018.07.006}.

\bibitem[Maurer et al.(2016)]{maurer2016}
Andreas Maurer, Massimiliano Pontil, and Bernardino Romera-Paredes.
\newblock The benefit of multitask representation learning.
\newblock \emph{Journal of Machine Learning Research}, 17(81):1--32, 2016.
\newblock \url{https://www.jmlr.org/papers/v17/15-242.html}.

\bibitem[Sch\"olkopf et al.(2021)]{scholkopf2021}
Bernhard Sch\"olkopf, Francesco Locatello, Stefan Bauer, Nan Rosemary Ke,
Nal Kalchbrenner, Anirudh Goyal, and Yoshua Bengio.
\newblock Toward causal representation learning.
\newblock \emph{Proceedings of the IEEE}, 109(5):612--634, 2021.
\newblock \url{https://doi.org/10.1109/JPROC.2021.3058954}.

\bibitem[Tripuraneni et al.(2020)]{tripuraneni2020}
Nilesh Tripuraneni, Michael I. Jordan, and Chi Jin.
\newblock On the theory of transfer learning: The importance of task diversity.
\newblock In \emph{Advances in Neural Information Processing Systems 33}, 2020.
\newblock \url{https://proceedings.neurips.cc/paper/2020/hash/59587bffec1c7846f3e34230141556ae-Abstract.html}.

\bibitem[Tripuraneni et al.(2021)]{tripuraneni2021}
Nilesh Tripuraneni, Chi Jin, and Michael I. Jordan.
\newblock Provable meta-learning of linear representations.
\newblock In \emph{Proceedings of the 38th International Conference on
Machine Learning}, pages 10434--10443, 2021.
\newblock \url{https://proceedings.mlr.press/v139/tripuraneni21a.html}.

\bibitem[von Luxburg(2007)]{vonluxburg2007}
Ulrike von Luxburg.
\newblock A tutorial on spectral clustering.
\newblock \emph{Statistics and Computing}, 17:395--416, 2007.
\newblock \url{https://doi.org/10.1007/s11222-007-9033-z}.

\end{thebibliography}

\end{document}
